\section{Presentacion convertida en practica: Rendimiento termico en procesadores}

Esta seccion nace de una presentacion existente y se incorpora al manual para que el material no permanezca aislado. La presentacion deja de ser un apoyo oral y se transforma en una ruta de estudio: problema, variables, datos, analisis, decision y actividad.

\begin{idea}
Equilibrar temperatura y rendimiento en computacion intensiva mediante diseno factorial y respuestas multiples.
\end{idea}

\subsection{Evidencia didactica extraida}

\begin{itemize}
\item Aplicación de un Diseño Factorial 3\textasciicircum{}3 y MANOVA para la Optimización de Energía y Ventilación.
\item El ANOVA para la variable Puntos muestra que los Hilos tienen un valor F inmenso (15,732), siendo el factor dominante.
\item Análisis Térmico: ANOVA y Efectividad del Ventilador.
\item Para la Temperatura , el ANOVA revela algo distinto: el factor Energía NO es significativo (p=0.20, vean la tabla), lo que significa que el CPU se calienta igual en modo 'Ahorro' o 'Alto'.
\end{itemize}

\subsection{Como entra al manual}

El tema se integra al capitulo relacionado con \textbf{manova factorial}. Si el estudiante solo observa la presentacion, ve una secuencia visual. En el manual, esa misma secuencia se expande con la pregunta de ingenieria, la razon del metodo, la interpretacion y los limites de la conclusion.

\subsection{Practica asociada}

La practica generada pide al estudiante transformar la presentacion en una decision: definir variables, organizar datos, graficar, aplicar el metodo y redactar una conclusion prudente. Si no existia practica, el sistema crea una primera version editable. Si existia una practica relacionada, esta se usa para enriquecer la narrativa y no como bloque aislado.

\subsection{Actividad de evaluacion}

Explique que parte de la presentacion corresponde a contexto, cual corresponde a diseno, cual corresponde a evidencia y cual corresponde a decision. Si falta alguno de estos elementos, proponga como completarlo.
